\documentclass[12pt]{article}
\usepackage[utf8]{inputenc}
\usepackage[T1]{fontenc}
\usepackage{hyperref}
\title{Ija-Zuba Language Resources for K-12 Curriculum Development}
\author{Compiled by Claude Sonnet 5 (claude-sonnet-5) for Andr\'as Kornai}
\date{}
\begin{document}
\maketitle

\section{Introduction}
The code vki (Ija-Zuba) was retired and split, effective 2021-01-15, into Koro Nulu [vkn] and Koro Zuba [vkz]. The two share very little vocabulary (about 7\% lexical similarity), and Koro Zuba appears to be Nupoid. The Glottolog entry for the retired code is \url{https://glottolog.org/resource/languoid/id/ijaz1238}. The task description gives the family as ``Bookkeeping'', which is Glottolog's label for retired codes and not a real language family. Curriculum work should target vkn and vkz separately. Search found no school materials for either.

\section{Primary Educational Resources}
I found no primers, textbooks or school programmes for Ija-Zuba, Koro Nulu or Koro Zuba.

\section{Linguistic and Reference Materials}
\begin{itemize}
\item Koro Zuba language (Wikipedia overview): \url{https://en.wikipedia.org/wiki/Koro_Zuba_language}
\item South Plateau languages (Wikipedia overview): \url{https://en.wikipedia.org/wiki/South_Plateau_languages}
\item Glottolog entry for the retired Ija-Zuba code: \url{https://glottolog.org/resource/languoid/id/ijaz1238}
\end{itemize}
The search summary mentions a 2021 sociolinguistic profile by Decker et al. in the \emph{Journal of Language Survey Reports}. The search did not return a URL for it, so it is not listed as a resource.

\section{Religious and Community Resources}
\begin{itemize}
\item Joshua Project profile, Koro Zuba in Nigeria: \url{https://joshuaproject.net/index.php/people_groups/20656}
\end{itemize}

\section{Digital and Audio Resources}
\begin{itemize}
\item Global Recordings Network, Koro Ija page: \url{https://globalrecordings.net/en/language/vki}
\end{itemize}
I did not verify what recordings this page actually offers.

\section{Curriculum Development Context}
Koro Zuba has an estimated 10,000--12,000 native speakers (2020) in about eighteen communities in Niger State and the FCT, Nigeria. Koro Nulu (Koro Ija) is reported as a single village near Lambata. Very little published material was found, so curriculum work would need new documentation. Next steps are to follow the vkn and vkz codes and the Decker et al. profile.
\end{document}